% TOP_PROBLEM: {"id": 3271, "title": "Arc-disjoint directed cycles in regular directed graphs", "category_id": 3, "category": {"id": 3, "name": "graph_theory", "display_name": "Graph Theory", "description": "Problems involving graphs, networks, and their properties.", "slug": "graph-theory", "order_index": 3, "created_at": "2026-07-31T15:26:25.671Z"}, "status": "open", "problem_number": "OPG-49573", "set_id": 12}
% FIRST_POSTED: 2026-10-01
% ENTRY_KIND: statement_only
% UnsolvedMath ID 3271; source catalogue code OPG-49573
% !TeX program = lualatex
% Definition and sources only; this file is not an LLM solution attempt.
\documentclass[11pt,a4paper]{article}
\usepackage[margin=25mm]{geometry}
\usepackage{fontspec}
\setmainfont{lmroman10-regular.otf}[BoldFont=lmroman10-bold.otf,
 ItalicFont=lmroman10-italic.otf,BoldItalicFont=lmroman10-bolditalic.otf]
\usepackage{amsmath,amssymb,mathtools}
\usepackage{unicode-math}
\setmathfont{latinmodern-math.otf}
\AtBeginDocument{\let\setminus\smallsetminus}
\usepackage{xurl,hyperref}
\hypersetup{colorlinks=true,urlcolor=blue}
\setcounter{secnumdepth}{0}
\setlength{\parindent}{0pt}
\setlength{\parskip}{5pt}
\setlength{\emergencystretch}{3em}
\begin{document}
\section*{Arc-disjoint directed cycles in regular directed graphs}\label{3271}
{\small UnsolvedMath ID 3271 · OPG-49573 · Graph Theory · OpenGarden}
\subsection{Definitions and mathematical statement}
Fix an integer $k\ge1$. A nonempty finite $k$-regular digraph $D=(V,A)$ here has
$d^+(v)=d^-(v)=k$ at every vertex. Assume it is loopless and has no
parallel arcs in the same direction. Opposite arcs are permitted and
are distinct members of $A$.

A directed cycle consists of distinct vertices joined cyclically by
arcs oriented in the direction of traversal. Its length is its number
of arcs; in particular an opposite pair of arcs is a cycle of length
two. A collection of cycles is arc-disjoint if no arc belongs to two
members. Vertices may be shared by many cycles.

The conjecture asks whether every such $D$ contains at least
\[
                             \binom{k+1}{2}=\frac{k(k+1)}2
\]
pairwise arc-disjoint directed cycles. It is enough to find exactly
that many, since extra cycles can be discarded. There is no requirement
that their union span all vertices or exhaust all arcs.

The coefficient is sharp for the complete symmetric digraph on $k+1$
vertices: it has exactly $k(k+1)$ arcs and every cycle uses at least two,
so an arc-disjoint collection contains at most $k(k+1)/2$ cycles. Taking
one digon for each unordered vertex pair attains that count. The
regularity convention is important: ``$k$-regular'' means both indegree
and outdegree equal $k$, not merely total degree $k$. The proposed lower
bound depends on this local degree, independently of the total order
of the digraph.
\subsection{Short English statement}
Must a digraph with indegree and outdegree k at every vertex contain at least k(k+1)/2 arc-disjoint directed cycles?
\subsection{Sources}
\begin{itemize}
\item[S1] ULAM.AI and the UnsolvedMath Contributors, supplied problems.json, record 3271 (OPG-49573), OpenGarden collection. Catalogue identity and source transcription; CC BY 4.0.
\url{https://huggingface.co/datasets/ulamai/UnsolvedMath}.
\item[S2] Open Problem Garden, Arc-disjoint directed cycles in regular directed graphs. Original problem and discussion; the mathematical formulation below is expanded from the supplied source record.
\url{http://www.openproblemgarden.org/op/arc_disjoint_directed_cycles_in_regular_directed_graphs}.
\end{itemize}
\end{document}
